% GENERATED FILE. Edit canonical lessons, then run npm run generate:books.
% canonical-source-hash: fnv1a64:326b2032a3644ce2
% canonical-lessons: MR-C265-parvangi, MR-C265-bandi, MR-C265-suraksa, MR-C265-salla, MR-C265-vacan

\chapter{Permission, Ban, Safety, Advice, Promise}
\label{ch:mr-permission-ban-safety-advice-promise}
\hlchaptermodality{265}

\begin{chapteropening}
\textbf{By the end of this chapter:} \emph{I can say permission, a ban, safety, advice and a promise.}

Complete the last lesson of chapter 265: I can say permission, a ban, safety, advice and a promise.
\end{chapteropening}

\section[\texorpdfstring{\mr{परवानगी} (parvāngī)}{परवानगी (parvāngī)}]{\mr{परवानगी} (parvāngī) — permission}

\label{lesson:MR-C265-parvangi}



\begin{quote}
Before the new one: say the Marathi for a father-in-law, then the Marathi for a son-in-law.
\end{quote}

\subsection*{You'll want to know: \mr{परवानगी}}
\textbf{\mr{परवानगी}} — \emph{parvāngī} — \textquotedblleft{}permission\textquotedblright{}.

\begin{cousinweb}
\textbf{\mr{परवानगी}} (\emph{parvāngī}) is a loan, usually traced to Persian \emph{parwānagī}, \textquotedblleft{}permission\textquotedblright{}, built on \emph{parwāna}, \textquotedblleft{}an order, a licence\textquotedblright{}.
\end{cousinweb}

\begin{grammarlens}[title={What you've built}]
One more word for this chapter.
\end{grammarlens}

\subsection*{Your turn}
\begin{itemize}
\raggedright
  \item \emph{Say it:} \emph{parvāngī}
  \item \emph{Say it:} \emph{parvāngī}, once more
  \item \emph{Say it:} say \emph{parvāngī}
  \item \emph{Recall:} say the Marathi for a father-in-law, then the Marathi for a son-in-law, then say \emph{parvāngī} again
\end{itemize}

\begin{tcolorbox}[breakable,colback=teal!4,colframe=teal!35!black,title={Before you move on}]
What does \mr{परवानगी} mean? (\textquotedblleft{}Permission\textquotedblright{}.) Say it once more.
\end{tcolorbox}
\section[\texorpdfstring{\mr{बंदी} (bandī)}{बंदी (bandī)}]{\mr{बंदी} (bandī) — a ban}

\label{lesson:MR-C265-bandi}



\begin{quote}
Before the new one: say the Marathi for a son-in-law, then the Marathi for permission.
\end{quote}

\subsection*{You'll want to know: \mr{बंदी}}
\textbf{\mr{बंदी}} — \emph{bandī} — \textquotedblleft{}a ban\textquotedblright{}.

\begin{cousinweb}
\textbf{\mr{बंदी}} (\emph{bandī}), \textquotedblleft{}a ban\textquotedblright{}, is built on \emph{band}, \textquotedblleft{}shut, closed\textquotedblright{}, a word usually traced to Persian \emph{band}.
\end{cousinweb}

\begin{grammarlens}[title={What you've built}]
One more word for this chapter.
\end{grammarlens}

\subsection*{Your turn}
\begin{itemize}
\raggedright
  \item \emph{Say it:} \emph{bandī}
  \item \emph{Say it:} \emph{bandī}, once more
  \item \emph{Say it:} say \emph{bandī}
  \item \emph{Recall:} say the Marathi for a son-in-law, then the Marathi for permission, then say \emph{bandī} again
\end{itemize}

\begin{tcolorbox}[breakable,colback=teal!4,colframe=teal!35!black,title={Before you move on}]
What does \mr{बंदी} mean? (\textquotedblleft{}A ban\textquotedblright{}.) Say it once more.
\end{tcolorbox}
\section[\texorpdfstring{\mr{सुरक्षा} (surakṣā)}{सुरक्षा (surakṣā)}]{\mr{सुरक्षा} (surakṣā) — safety, security}

\label{lesson:MR-C265-suraksa}



\begin{quote}
Before the new one: say the Marathi for permission, then the Marathi for a ban.
\end{quote}

\subsection*{You'll want to know: \mr{सुरक्षा}}
\textbf{\mr{सुरक्षा}} — \emph{surakṣā} — \textquotedblleft{}safety, security\textquotedblright{}.

\begin{cousinweb}
\textbf{\mr{सुरक्षा}} (\emph{surakṣā}) is Sanskrit \emph{su-}, \textquotedblleft{}well\textquotedblright{}, and \emph{rakṣā}, \textquotedblleft{}protection\textquotedblright{}, a tatsama.
\end{cousinweb}

\begin{grammarlens}[title={What you've built}]
One more word for this chapter.
\end{grammarlens}

\subsection*{Your turn}
\begin{itemize}
\raggedright
  \item \emph{Say it:} \emph{surakṣā}
  \item \emph{Say it:} \emph{surakṣā}, once more
  \item \emph{Say it:} say \emph{surakṣā}
  \item \emph{Recall:} say the Marathi for permission, then the Marathi for a ban, then say \emph{surakṣā} again
\end{itemize}

\begin{tcolorbox}[breakable,colback=teal!4,colframe=teal!35!black,title={Before you move on}]
What does \mr{सुरक्षा} mean? (\textquotedblleft{}Safety, security\textquotedblright{}.) Say it once more.
\end{tcolorbox}
\section[\texorpdfstring{\mr{सल्ला} (sallā)}{सल्ला (sallā)}]{\mr{सल्ला} (sallā) — advice}

\label{lesson:MR-C265-salla}



\begin{quote}
Before the new one: say the Marathi for a ban, then the Marathi for safety.
\end{quote}

\subsection*{You'll want to know: \mr{सल्ला}}
\textbf{\mr{सल्ला}} — \emph{sallā} — \textquotedblleft{}advice\textquotedblright{}.

\begin{cousinweb}
\textbf{\mr{सल्ला}} (\emph{sallā}) is a loan, usually traced through Persian to Arabic \emph{ṣalāḥ}, \textquotedblleft{}soundness, what is right\textquotedblright{}, whose root is \emph{ṣ-l-ḥ} (Lane).
\end{cousinweb}

\begin{grammarlens}[title={What you've built}]
One more word for this chapter.
\end{grammarlens}

\subsection*{Your turn}
\begin{itemize}
\raggedright
  \item \emph{Say it:} \emph{sallā}
  \item \emph{Say it:} \emph{sallā}, once more
  \item \emph{Say it:} say \emph{sallā}
  \item \emph{Recall:} say the Marathi for a ban, then the Marathi for safety, then say \emph{sallā} again
\end{itemize}

\begin{tcolorbox}[breakable,colback=teal!4,colframe=teal!35!black,title={Before you move on}]
What does \mr{सल्ला} mean? (\textquotedblleft{}Advice\textquotedblright{}.) Say it once more.
\end{tcolorbox}
\section[\texorpdfstring{\mr{वचन} (vacan)}{वचन (vacan)}]{\mr{वचन} (vacan) — a promise}

\label{lesson:MR-C265-vacan}



\begin{quote}
Before the new one: say the Marathi for safety, then the Marathi for advice.
\end{quote}

\subsection*{You'll want to know: \mr{वचन}}
\textbf{\mr{वचन}} — \emph{vacan} — \textquotedblleft{}a promise\textquotedblright{}.

\begin{cousinweb}
\textbf{\mr{वचन}} (\emph{vacan}) is Sanskrit \emph{vacana}, \textquotedblleft{}speaking; statement\textquotedblright{}, taken over whole as a learned loan, a tatsama (CDIAL 11199). It is built on \emph{vac}, \textquotedblleft{}to speak\textquotedblright{}.
\end{cousinweb}

\begin{grammarlens}[title={What you've built}]
One more word for this chapter.
\end{grammarlens}

\subsection*{Your turn}
\begin{itemize}
\raggedright
  \item \emph{Say it:} \emph{vacan}
  \item \emph{Say it:} \emph{vacan}, once more
  \item \emph{Say it:} say \emph{vacan}
  \item \emph{Recall:} say the Marathi for safety, then the Marathi for advice, then say \emph{vacan} again
\end{itemize}

\begin{tcolorbox}[breakable,colback=teal!4,colframe=teal!35!black,title={Before you move on}]
What does \mr{वचन} mean? (\textquotedblleft{}A promise\textquotedblright{}.) Say it once more.
\end{tcolorbox}
